\documentclass[11pt]{article}
\usepackage{amsmath, amssymb, geometry, hyperref}
\geometry{a4paper, margin=1in}

\title{Work Stream 3: Pre-Singularity Journey \& Turbulence Modeling}
\author{MechanicaFluidorum Program \& Community Contributors (esp. \texttt{u/babainottawa})}
\date{September 2026}

\begin{document}
\maketitle

\section{Introduction}
As noted by \texttt{u/babainottawa}, while the singularity itself is physically unreachable, the \emph{journey} of the singularity contains valuable information. The extreme vortex stretching observed in the regime where $Ma < 0.3$ offers a unique mathematical testbed for turbulence models.

\section{Subgrid-Scale Stress and Singularity Trajectories}
In Large Eddy Simulation (LES), the subgrid-scale (SGS) stress tensor $\tau_{ij}$ is often modeled using the Smagorinsky eddy viscosity model. By taking the exact mathematical trajectory of the OpenAI blowup construction, we can a priori test how standard SGS models respond to extreme energy cascades.

The oscillating wave packets constructed in the proof transfer energy to the core vortex in a highly structured, anisotropic manner. Analyzing this process provides a benchmark for dynamic models (like the Germano identity) under transient, non-equilibrium turbulence conditions.

\section{Acknowledgments}
We thank \texttt{u/babainottawa} for shifting the focus from the mathematical endpoint to the physical journey. \textit{Rendons \`a C\'esar ce qui appartient \`a C\'esar.}

\end{document}
